\documentclass[11pt,a4paper]{article}
\usepackage[T1]{fontenc}
\usepackage{mathptmx}
\usepackage[margin=25mm]{geometry}
\usepackage{booktabs}
\usepackage{array}
\usepackage{amsmath}
\usepackage{microtype}
\usepackage{titlesec}
\usepackage[hidelinks]{hyperref}
\usepackage{fancyhdr}
\titleformat{\section}{\large\bfseries}{\thesection}{0.6em}{}
\titleformat{\subsection}{\bfseries}{\thesubsection}{0.6em}{}
\setlength{\parskip}{0.5em}
\setlength{\parindent}{0pt}
\pagestyle{fancy}\fancyhf{}
\renewcommand{\headrulewidth}{0.4pt}
\newcommand{\inr}[1]{INR~#1}

\fancyhead[L]{RescueSwarm --- Track B: Avionics and Communications}
\fancyhead[R]{\thepage}
\begin{document}

\begin{center}
{\LARGE\bfseries Track B --- Avionics and Communications}\\[3mm]
{\large RescueSwarm: work package and funding request}\\[2mm]
NIDAR 2026--27 $\cdot$ Track 1 $\cdot$ Mission 1\\[1mm]
\textbf{Ask: \inr{280,963}} $\cdot$ 1--2 students $\cdot$ Flight controller, companion computer, GNSS/RTK, mesh, safety link, video
\end{center}
\vspace{2mm}\hrule\vspace{4mm}


This track owns everything electrical between the airframe and the autonomy.
It carries the second-largest ask in the programme, and it carries the single
component whose specification is currently in doubt.

\section{Scope and ownership}

This track owns \textbf{flight controller, companion computer, gnss/rtk, mesh, safety link, video}. It is one of four delivery tracks defined in
\texttt{docs/development-plan.md}~\S1.3; the integrated system argument lives in
the master proposal, \texttt{docs/proposal/rescueswarm-proposal.pdf}, which this
document does not restate.

\section{Compute, and why it is three parts rather than one}

The stack is a flight controller, a companion computer and an accelerator, and
the division is not arbitrary. The accelerator is an M.2 class device whose only
interface is \textbf{PCIe}; it is a neural processing unit, not a computer, and
it cannot boot, hold a filesystem or terminate a camera link. The camera is a
CSI module producing a Bayer-mosaic stream that is not an image until something
demosaics and scales it. The companion computer is what both plug into.

\textbf{This has a hard consequence for procurement.} The accelerator requires a
host that exposes PCIe. That rules out the previous-generation single-board
computer entirely --- its PCIe lane is committed internally to a USB controller
and is never broken out. There is no adapter and no workaround. The host
generation is a dependency of the accelerator, not a performance preference,
and any cost pass that treats it as one produces a bill of materials that cannot
be assembled.

\section{Positioning: the item this track must resolve first}

Geolocation accuracy governs \textbf{125 of 200} available geotagging points, and
the error budget is unambiguous about what drives it:

\begin{center}
\begin{tabular}{@{}lr@{}}
\toprule
\textbf{Configuration} & \textbf{RSS error} \\
\midrule
Standard GNSS, uncalibrated boresight & 4.57\,m \\
Standard GNSS, calibrated boresight & 3.88\,m \\
RTK + dual-antenna heading + calibration & 0.75\,m \\
\quad + 20-frame fusion & 0.66\,m \\
\bottomrule
\end{tabular}
\end{center}

The step from 3.88\,m to 0.75\,m is RTK. Nothing else in the budget moves the
number comparably, which is why the receiver and its base station survive every
cost-reduction pass.

\textbf{The open item.} The receiver currently carried against this line is
specified by its supplier as supporting SBAS correction services with a CEP
below 1.5\,m. That is assisted GNSS, not RTK, and it lands the system in the
3.88\,m row rather than the 0.75\,m row. Resolving this --- by confirming the
part, substituting a genuine RTK receiver, or re-scoring the geotagging
expectation --- is this track's first P1 action, because 125 points depend on
which row is true.

\section{The radio link, and a deferral that depends on a configuration}

The sub-GHz safety radio is deferred on the strength of a specific claim: that
the primary radio link runs in a firmware mode carrying both control and MAVLink
telemetry on one link and one autopilot serial port. The older transparent-serial
mode \emph{consumes} that link --- configured that way the aircraft has telemetry
and no control, and the deferred radio is required again.

\textbf{This is not a preference, it is the precondition of a deferral}, and it
must be verified on the bench in P2 rather than assumed.

The mesh link budget is thin and this track should say so plainly: 5.8\,GHz
carries \textbf{1.7\,dB} of margin at the 600\,m design range against a 15\,dB
target for a mobile airborne link. The mitigation is architectural rather than
financial --- keep the ground antenna elevated, stream one switched video feed
rather than three, and accept that video degrades before command does. Total
offered load is 2.5\,Mbps against a 20\,Mbps channel, so the constraint is
margin, not bandwidth.

\section{Failsafe configuration}

The battery failsafe thresholds are derived quantities, not chosen ones:
\texttt{BATT\_LOW\_VOLT} 18.48\,V and \texttt{BATT\_CRT\_VOLT} 17.18\,V follow
from the cell's resting curve and a pack resistance of
\texttt{BATT\_RESISTANCE}~=~0.0400\,$\Omega$. A previous configuration set the
sag-compensated voltage source while leaving pack resistance unset, so the
compensation did nothing and a loaded voltage was compared against a threshold
taken from the resting curve. That class of defect --- a correct-looking
configuration that nothing actually applies --- is what this track's bench
verification exists to catch.

\section{Funding request}

Every figure below is generated from
\texttt{docs/proposal/figures/track\_budget.py}, which partitions the master
competition budget. The four track asks plus the shared pool reconcile to the
master figure exactly; that reconciliation is asserted in code and fails the
build if it ever stops holding. \textbf{K} marks a line held at professional
grade on purpose.


\begin{center}
\begin{tabular}{@{}p{88mm}cr@{}}
\toprule
\textbf{Item} & & \textbf{\inr{}} \\
\midrule
Flight controller (1 x 22,600 x 3 ac) & \textbf{K} & 67,800 \\
GNSS RTK primary (1 x 18,000 x 3 ac) & \textbf{K} & 54,000 \\
RC rx, storage, cooling, mounts (1 x 8,500 x 3 ac) &  & 25,500 \\
Companion computer (1 x 8,000 x 3 ac) & \textbf{K} & 24,000 \\
Mesh node + antennas (1 x 6,000 x 3 ac) &  & 18,000 \\
RTK base receiver &  & 18,000 \\
Safety-pilot transmitter &  & 8,000 \\
Survey tripod + tribrach &  & 4,000 \\
\midrule
Subtotal & & 219,300 \\
Customs duty and freight & & 10,217 \\
GST & & 14,799 \\
Contingency at 15\% & & 36,647 \\
\midrule
\textbf{TRACK B ASK} & & \textbf{280,963} \\
\bottomrule
\end{tabular}
\end{center}


\textbf{Deferred or institutionally held, and therefore not requested:}
Spare GNSS, Spare compute module, Spare flight controller, VEGA co-processor (1 x 0 x 3 ac), WPC / ETA licensing. These remain capabilities the track requires; they are simply not being
bought. Where a deferral carries a technical consequence it is stated in the
risk table below rather than left implicit.

\section{Deliverables by phase}

\begin{center}
\begin{tabular}{@{}llp{88mm}@{}}
\toprule
\textbf{Phase} & \textbf{Dates} & \textbf{This track delivers} \\
\midrule
\textbf{P1} & 24 Aug -- 13 Sep & GNSS specification resolved. Long-lead orders: compute, camera, radios, GNSS. Interface control documents agreed. \\
\textbf{P2} & 14 Sep -- 4 Oct & Radio mode verified on the bench. Failsafe thresholds validated against a real pack discharge. \\
\textbf{P3} & 5 Oct -- 18 Oct & Link range verified in the field. RTK fix demonstrated. \\
\textbf{P5} & 9 Nov -- 22 Nov & All links exercised under mission load. \\
\bottomrule
\end{tabular}
\end{center}

\section{Risks owned by this track}

\begin{center}
\begin{tabular}{@{}p{50mm}lp{70mm}@{}}
\toprule
\textbf{Risk} & \textbf{Sev.} & \textbf{Mitigation} \\
\midrule
Primary GNSS is not RTK-capable & High & Resolve at P1. If confirmed, either substitute a true RTK receiver or restate the geotagging expectation. 125 points ride on this. \\
Radio mode does not carry control and telemetry together & High & Verify on the bench at P2. If it does not, the deferred sub-GHz safety radio must be restored, which is a budget event. \\
5.8\,GHz margin is 1.7\,dB against a 15\,dB target & Medium & Elevated ground antenna, single switched video feed. Command link is on a separate budget with 12.2\,dB. \\
\bottomrule
\end{tabular}
\end{center}

\section{Interfaces to other tracks}

\begin{center}
\begin{tabular}{@{}lp{110mm}@{}}
\toprule
\textbf{With} & \textbf{Dependency} \\
\midrule
Track A & Depends on regulated power and on vibration-isolated mounting. \\
Track C & Supplies the companion computer the autonomy runs on, and the MAVLink routing it commands through. \\
Track D & Supplies the camera interface, the accelerator, and the time and attitude reference every geotag is computed against. \\
\bottomrule
\end{tabular}
\end{center}

Interfaces are where student programmes fail, because each side assumes the
other owns the gap. Every row above is a two-way commitment and should be
recorded as an interface control document at P1.


\vfill
\hrule\vspace{2mm}
{\small Generated by \texttt{tools/proposal/build\_track\_proposals.py} from
\texttt{docs/proposal/figures/track\_budget.py}. Financial figures are not
maintained in this document and must not be edited here: change the master
budget and regenerate. Technical claims cite the model or document that
produced them.}
\end{document}
